\label{sec:discussion}

\subsection{Comparison to Commission-Drawn Maps}

The most natural comparators for \textsc{bisect} algorithmic maps are commission-drawn
and court-drawn maps, which are produced by humans but under instructions to apply
neutral criteria without partisan intent. Unfortunately, direct comparison is limited
by data availability: NC's 2022 map was drawn by the legislature and partially modified
by courts; WI's 2022 map was drawn by the legislature and partially modified by the
state Supreme Court; TX's 2022 map was drawn by the legislature and upheld with minor
modifications after litigation.

For NC, we can compare to the court-drawn remedial map imposed by the NC Supreme Court
in early 2022, which was subsequently set aside by the conservative majority reinstated
after the 2022 election. The court-drawn NC map produced approximately 3 open seats
and a safe-seat fraction of approximately 0.50---closer to the \textsc{bisect} outcome
(4 open seats, 0.43 safe-seat fraction) than to the enacted outcome (2 open seats, 0.64
safe-seat fraction). This suggests that neutral human redistricting produces open-seat
and safe-seat outcomes more similar to algorithmic redistricting than to enacted plans.

Colorado provides the best available commission comparison: the Independent
Redistricting Commission's 2022 congressional map produced 3 open seats out of 8
districts, a fraction of 0.375---comparable to the algorithmic open-seat rates reported
in Table~\ref{tab:open-results} for WI. The Colorado commission was explicitly prohibited
from considering incumbent addresses in its deliberations, making it the closest
available human-drawn analogue to the algorithmic approach.

\subsection{Implications for Electoral Accountability}

The open-seat analysis has a direct implication for electoral accountability that
extends beyond the individual retirement decision. A redistricting plan's open-seat
count determines, in part, the rate at which seats change hands between parties across
election cycles: open seats are more likely to change party than incumbent seats, and
more open seats mean more competitive elections and more responsive legislative
delegations.

From this perspective, the enacted maps' low open-seat counts represent a systematic
reduction in electoral accountability. Incumbents who are shielded from competitive
elections by engineered home districts face fewer incentives to maintain alignment
with their constituents' preferences and are less vulnerable to party-level accountability
through wave elections. The \textsc{bisect} comparison demonstrates that a geographically
compact map produces substantially more open seats---not as a goal, but as a consequence
of geographic optimisation that does not protect incumbent residences.

\subsection{No Strategic Intent}

A critical feature of the open-seat analysis is that the algorithmic open seats are
not the product of any strategic intent to increase electoral competition. The algorithm
does not know that creating a district without an incumbent home tract produces an open
seat; it simply draws compact, population-balanced districts without regard to where
incumbents live. The open seats that result are geometric accidents: the compact district
boundary happens to exclude all incumbent home tracts from a particular district,
producing an open seat as a by-product of geographic optimisation.

This distinguishes algorithmic redistricting from proposals to deliberately maximise
competitiveness as a redistricting criterion. Proposals to mandate competitive districts
(e.g., requiring at least $X\%$ of districts to have partisan lean within $Y$ percentage
points) are constitutionally contested: some state courts have accepted competitiveness
as a criterion, others have rejected it. The \textsc{bisect} algorithm's indifference
to competition means that its maps cannot be challenged as pursuing an impermissible
competitiveness mandate---the competition is a consequence of compactness, not a target
of the algorithm.

\subsection{Retirement Deterrence vs.\ Forced Retirement}

The retirement incentive model distinguishes between retirement deterrence and forced
retirement. Retirement deterrence occurs when an incumbent faces a more competitive
district or a stronger primary opponent, increasing the probability of choosing to
retire rather than run. Forced retirement occurs when two incumbents are paired and one
is effectively unable to win the primary, producing a near-certain retirement.

Under \textsc{bisect} plans, both mechanisms are active. The pairing events documented
in I.1 produce forced retirements: one of the paired incumbents will almost certainly
leave Congress after the redistricting cycle, whether by losing the primary, relocating
to an adjacent open seat, or retiring. The partisan composition changes documented in
I.2 produce retirement deterrence: incumbents whose safe seats become competitive face
higher campaign costs and lower reelection probabilities, which translates into higher
retirement probabilities in the aggregate.

The combined effect, captured in the retirement incentive score, represents the full
scope of the algorithmic redistricting disruption to incumbency. Enacted maps
deliberately minimise both mechanisms: they separate incumbents to avoid pairing and
maintain safe-seat margins to prevent competitive deterrence. The \textsc{bisect}
algorithm, by being indifferent to both, allows both mechanisms to operate at their
geographic baseline rates.
